% Proposed replacement for \subsection*{Data sources} in the Overleaf main.tex (28 October 2026), checked against the
% Snow-polygon primary (primary_snow520_merge_dhsmics_gamma2_v12; the counts quoted are the same in v11). main.tex is not edited; paste by hand.
% New citation keys needed in ref.bib: igme (UN IGME all-cause estimates), wpp (UN World Population Prospects).
\subsection*{Data sources}

Complete birth histories were obtained from the Births Recodes of the Demographic and Health Surveys and Malaria Indicator Surveys (DHS/MIS) \cite{dhs} and from UNICEF Multiple Indicator Cluster Surveys (MICS) that recorded a complete birth history \cite{mics}. Birth dates, survival status at the time of the survey interview, and reported age at death were used to construct individual child trajectories. Ages were grouped in 7 age-bands: $<1$ month, 1--5 months, 6--11 months, 1--2 years, 2--3 years, 3--4 years and 4--5 years. An age band was eligible if the child entered it alive in the 60 months before the interview and in a calendar year between 2000 and 2024, and the whole band ended before the interview, irrespective of survival.

The units of analysis were the 520 malaria-endemic sub-national polygons defined by Snow et al.\ \cite{snow:17}. Each record was located by the smallest administrative unit reported in its survey: the survey region (defined by the survey's DHS boundary polygons or, for MICS, by the region mapping described below) or, where a survey region spanned several polygons, a finer location variable (for example state, province or district) matched by boundary geometry, polygon name or a reviewed alias. A record's unit was the polygon holding at least 90\% of the population of its location (Gridded Population of the World version 4 (GPWv4), 2020 \cite{gpwv4}), or otherwise the union of the polygons each holding at least 5\%. Units with more than half of their population in malaria-free or island polygons were excluded. The included surveys contributed 1,452 survey units (a survey's records in one unit), covering 465 polygons; 215 survey units were unions of two or more polygons.

Thirteen area-level covariates were computed for each survey unit from the survey microdata, using the DHS indicator definitions and survey weights. From the DHS/MIS Births Recodes: the mean maternal age at first live birth, mean maternal education in completed years, mean household wealth quintile and percentage of urban residence among mothers with a live birth in the 60 months before interview; coverage of three doses of a DTP-containing vaccine (DTP3) and of a measles-containing vaccine among living children aged 12--23 months; the percentage of births in the 60 months before interview delivered in a health facility; and the percentage of non-first births with a preceding birth interval shorter than 24 months. From the household and height-and-weight recodes: the percentage of households with an improved drinking-water source, an improved sanitation facility and electricity, and the prevalence of wasting and stunting (below $-2$ SD of the WHO Child Growth Standards) among de facto children under five years. National values computed in the same way were checked against the published DHS estimates. Missing DTP3 or measles coverage was replaced by the WHO/UNICEF national estimate (WUENIC) for the survey year, and any covariate still missing in a survey unit took the mean of the survey's other units, weighted by their number of records.

Of the 93 MICS surveys from sub-Saharan Africa, 47 (28 countries) recorded a complete birth history; the other 46 recorded only summary birth histories and thus cannot be used in the age-band analysis. Each MICS birth history was converted to the DHS Births Recode layout and processed with the same age-band construction and eligibility rules as the DHS/MIS data. MICS does not publish region boundaries, so the 366 MICS region labels were matched, by exact name (291) or a reviewed correspondence table (75), to the regions of a DHS survey in the same country, whose boundaries are published; where no DHS boundaries exist, they were matched to admin-1 polygons (Guinea-Bissau and the Central African Republic) or directly to the Snow et al.\ polygons (South Sudan). The 13 covariates were computed from the MICS files with the same definitions; education, recorded as level and grade, was converted to completed years, and facility delivery refers to the last birth in the two years before interview.

\emph{P.\ falciparum} prevalence estimates were taken from the Malaria Atlas Project (MAP) annual modelled surfaces of the parasite rate standardised to ages 2 to 10 years (\textit{Pf}PR$_{2-10}$) for 2000--2024 \cite{map,map2025}. Each polygon's value was the mean of the surface over the polygon, weighted by GPWv4 2020 population counts (population density $\times$ cell area) on the MAP grid; the value for a union of polygons was the population-weighted mean of its polygons \cite{gpwv4}.

Each age band was assigned the \textit{Pf}PR$_{2-10}$ of its unit in the calendar year of entry into the band. For national burden calculations, country-level prevalence was calculated using population-weighted means, with weights derived from the GPWv4 population density surface for 2020 and grid-cell areas \cite{gpwv4}.

National malaria-attributable mortality was calculated from UN Inter-agency Group for Child Mortality Estimation (UN IGME) estimates of all-cause deaths at $<1$ month, 1--11 months and 1--4 years \cite{igme}, with the UN World Population Prospects population under 5 as the child-years at risk \cite{wpp}. Each UN IGME age group was split across the model's age bands in proportion to the deaths observed in the survey sample in that group. In a sensitivity analysis, all-cause deaths and person-years by age were taken from the IHME Global Burden of Disease estimates \cite{ihme}. Comparator estimates of under-five malaria deaths were obtained from IHME and from the CA-CODE 2026 series disseminated through the UN IGME portal \cite{Perin:26,Villavicencio:24}. We also used IHME state-level estimates of all-cause deaths by age in Nigeria.

National child HIV incidence estimates for ages 0--14 years, expressed as new infections per 1,000 uninfected population, were obtained from the UNAIDS 2025 estimates distributed by UNICEF. Corresponding adolescent incidence estimates for ages 15--19 years informed the imputation of missing child incidence; Liberia, which has no UNAIDS/UNICEF incidence series, was assigned the rate implied by UNAIDS counts of new infections in children aged 0--14 years. National GDP per capita, current health expenditure per capita (both in current US\$) and political stability and absence of violence were obtained from the World Bank World Development Indicators and Worldwide Governance Indicators \cite{wdi,wgi}. These four national covariates were matched to the country and calendar year of band entry; for years before South Sudan's independence in 2011, South Sudan took Sudan's values for GDP, health expenditure and political stability.
